\chapter{The Software Stack}
\label{ch:stack}

An analysis depends on several layers of software, from the operating system
to the scripts. Understanding what each layer does helps us choose tools and
keep the workflow reproducible.

\section{Layers and interfaces}

\begin{description}
  \item[Operating system] Linux is the path of least resistance; macOS is close
    and has an advantage for local models (\Cref{ch:hardware}). The shell
    listings in this book assume one of the two. On Windows they need WSL2, a
    Linux environment inside Windows, which an RStudio user has to install and
    learn before the first listing runs. The R listings run unchanged. macOS
    ships BSD versions of \code{grep}, \code{sed} and \code{find}, which lack
    some GNU options, and a command that fails there with ``invalid option''
    usually fails for that reason.
  \item[Languages and environments] R for ecological statistics, Python for
    machine learning; most projects use both. Pin packages with \code{renv}
    for R and \code{uv} or \code{conda} for Python.
  \item[Containers] Podman or Docker, for hard reproducibility cases and for
    fencing in agents. Podman runs without root privileges by default.
  \item[Editor, harness, version control] Positron, RStudio, VS Code or Neovim,
    plus an agent harness and Git with a forge.
  \item[Workflow engines] \code{targets} for R, Snakemake or Nextflow for
    multi-tool pipelines: they track dependencies and rerun only what changed.
\end{description}

There are four ways to talk to a system. A graphical interface (GUI) is easy to
explore and cannot be scripted. A text interface in the terminal (TUI) works
over a remote connection. A command line (CLI) can be scripted and logged. An
API does anything, if you write the code. For reproducibility the command line
wins: a command is a record of what you did, a series of clicks is not.
Harnesses run in the terminal, inside editors and in the cloud; this book
uses the terminal kind because a session there leaves a record that can be
read and replayed (\Cref{fig:layers}).

\begin{figure}[htbp]
\centering
\fig{ch08-layers}
\caption{The layers an analysis depends on, and the four ways to talk to a
system. The command line can be scripted and logged, which is why this book's
harnesses run there.}
\label{fig:layers}
\end{figure}

\section{Ten rules for code that lasts}

\begin{enumerate}
  \item Minimise dependencies; each is a maintenance and security liability.
  \item Pin exact versions in a lockfile.
  \item Keep everything textual under version control.
  \item One script, one job, numbered, with defined inputs and outputs.
  \item No absolute paths (\code{here::here()} in R, \code{pathlib} in Python).
  \item Set and record random seeds.
  \item Assert your assumptions: row counts, value ranges, join cardinality.
  \item Write the README as you go.
  \item Run everything from raw data to figures with one command.
  \item Evaluate a tool before depending on it, using SNIFF. Its criteria are \textbf{S}ource
    (who maintains it and why), \textbf{N}etwork (is it part of an ecosystem),
    \textbf{I}nteraction (are developers and users active), \textbf{F}it (use
    case and licence) and \textbf{F}ragility (its own dependencies)
    \parencite{nagraj2026sniff}.
\end{enumerate}

These practices also make an agent's work easier to check. A disorganised
project lets mistakes spread quickly, even when the agent appears to be
making progress. For biologists who want these skills from the start,
\textcite{allesina2019computing} give a chapter each to the Unix shell,
version control, writing good code, regular expressions, statistical
computing and relational databases.

\section{Local model tooling}

\textbf{Ollama} is the usual starting point: it downloads and runs models and
offers an OpenAI-compatible API on \code{localhost:11434}. \textbf{llama.cpp}
sits underneath much of this and gives finer control; \textbf{vLLM} serves many
users at once; \textbf{LM Studio} is the point-and-click option. On Linux, run
the runtime in a rootless container; \textcite{monshed2026ollamapodman} give a
working Podman recipe that also solves GPU access. \textbf{Open WebUI} provides
a familiar chat window over local and remote models. Interfaces with extensions
can rewrite what goes in and out: Lars Juhl Jensen's \emph{Seshat}, listed in
\Cref{app:tutorials}, adds PubMed results to a prompt before the model sees
it, a domain RAG layer as a plug-in. On Apple silicon, MLX, Apple's array
framework, is the native route; LM Studio and Ollama both offer MLX builds,
and the sessions in \Cref{ch:firstsession} ran on one (\Cref{app:stateofplay}).
One warning: Ollama also serves models hosted on its own servers through the
same command line, and a prompt sent to one of those leaves the machine.
Check that a model is local before trusting it with data
(\Cref{ch:localfirst}).

\begin{lstlisting}[style=shell,caption={A local coding model with a usable context}]
ollama pull MODEL                  # a current tag is in Appendix D
ollama run MODEL
  /set parameter num_ctx 65536     # the default is sized for chat, not agents
  /save MODEL-64k
ollama ps                          # CONTEXT must show the new value
\end{lstlisting}

Two settings cause most local-agent failures (\Cref{fig:localtooling}). The harness must be told that
the model supports tools (in OpenCode, \code{"tool\_call": true} in the model's
entry). The context must also be enlarged. The runtime's default is sized for
chat and depends on the machine's memory. An agent loop re-sends the whole
session each turn and needs tens of thousands of tokens; otherwise information
is cut and the agent appears unable to follow the task. Set it per model as
above, or
once for the server with \code{OLLAMA\_CONTEXT\_LENGTH} as
\Cref{ch:firstsession} does. \Cref{app:stateofplay} gives the current defaults
and the vendor's advice, and \Cref{ch:hardware} the memory a long context
costs. A local model of about twelve billion parameters is useful for small
edits and explanation, and well behind frontier models on work spanning many
files.

\begin{figure}[htbp]
\centering
\fig{ch08-local-tooling}
\caption{Local model tooling. Clients speak one OpenAI-compatible API, runtimes
serve the model, and llama.cpp sits underneath much of it. The context length
and the tool-calling flag cause most local-agent failures.}
\label{fig:localtooling}
\end{figure}

\section{Speech to text}

Speech to text is a different technology from the species-call classifiers
used in acoustic monitoring, such as BirdNET, and nothing here carries over to
them. \code{whisper.cpp} runs the Whisper transcription model locally. On
Apple silicon its default build runs on the GPU. An optional Core ML build
moves the encoder to the Neural Engine, which the project reports as more
than three times faster than running on the CPU alone, not than the GPU
build. Beware: Whisper invents text over silence, music or noise. In one
study, roughly 1 per cent of audio transcriptions contained entire phrases or
sentences that were never spoken, more often for speakers with long pauses
\parencite{koenecke2024whisper}.

Greek shows how this looks in practice. When \textcite{paraskevopoulos2024greekpodcast}
transcribed 800 hours of Greek podcasts, Whisper wrote a subtitle credit,
``Ypotitloi AUTHORWAVE'' (``AUTHORWAVE subtitles''), over music that voice
detection had missed. Removing those segments, the detected silences and passages
mixing Greek with English reduced the corpus by about 20 per cent. Such phantom
phrases recur often enough to be listed and filtered after transcription
\parencite{baranski2025whisper}. Filter silence first, remove the known phrases,
and never treat an automatic transcript as ground truth (\Cref{fig:speech}) --- least of all an
interview, where an invented sentence is attributed to a real person.

\begin{figure}[htbp]
\centering
\fig{ch08-speech}
\caption{Speech to text, with the checks around it. Whisper invents text over
silence and music, so filter before transcribing and remove the known phantom
phrases afterwards.}
\label{fig:speech}
\end{figure}
